\chapter{Dataset}
\label{ch:datasets}

This chapter describes the datasets used for training and evaluating the TRNet and MISRGRU architectures.
All experiments use the same synthetic dataset that was created by Komen and Tekpinar, which was used 
in the thesis~\cite{jelle2024} and subsequently in the RESURF paper~\cite{resurf2024}.
No modifications were made to the simulation pipeline or to the raw simulated data.

In this work we motion augmentated simulated data to train and evaluate the network capability on reconstructing motion.
Section~\ref{sec:motion_augment} describes how motion augmentation was achieved.

\section{Synthetic Data }
\label{sec:synthetic}

Fluorescence microscopy experiments are time-consuming and thus obtaining
ground-truth super-resolved images from real data is hard. For this reason, a simulation-based 
dataset was used where different conditions could be controlled such as blinking speed and 
background noise. 
The synthetic dataset provides paired low-resolution (LR)
fluorescence stacks of 100 frames and corresponding high-resolution (HR) SOFI targets.

\subsection{Simulated Data Conditions}
\label{sec:simulation_pipeline}

To build a dataset that generalises across different experimental conditions,
simulations were run under 12 distinct signal-to-noise settings formed by the
combination of three blinking-speed regimes and four photon-level fractions \cite{komen2024}:

\begin{itemize}
  \item \textbf{Blinking speeds}: fast, medium, and slow, corresponding to
        fluorophore on/off-times of approximately 600ms, 1200ms, and 2000ms respectively.
  \item \textbf{Photon levels}: 100\%, 80\%, 60\%, and 40\% of the nominal ion photon count,
        progressively reducing the peak signal intensity.
\end{itemize}


\section{Motion Augmentation}
\label{sec:motion_augment}

The synthetic dataset described above contains no sample motion: every frame in a
sequence depicts the same field of view, with fluorophores at fixed coordinates
emitting stochastic blinks.
This is appropriate for the batch MISRGRU and TRNet models, which output a single
SR image from a fixed set of $N$ frames.
The streaming MISRGRU (Section~\ref{sec:streaming_arch}), however, is designed to
produce a super-resolved image at \emph{every} time step, making it directly
relevant to live-cell microscopy where biologists monitor dynamic processes
in real time.

In live-cell imaging, sample drift on the order of tens of nanometres per second
is common and unavoidable.
Organelles such as vesicles and mitochondria can translate across the field of view
at 0.5--5\,\textmu m/s~\cite{tang2013fast}.
A model trained exclusively on static data has no incentive to track the current
field position: the correct answer is always the same single HR reconstruction. Our experiments showed staticaly trained 
models produced blurs where they summed all locations of fluorophores rather then reconstructed them at new location. \todo{ref experiment showing ghost blurs}  
Motion augmentation transforms each static training sample into a sequence with a
plausible slow drift trajectory and supplies a matching per-frame HR target,
forcing the model to localise the field of view at each time step.
Because paired moving ground-truth data cannot be obtained experimentally, all
motion is synthesised from the existing static dataset.
An earlier approach considered perturbing microtubule segment angles frame-by-frame
using an Ornstein-Uhlenbeck process; this was abandoned owing to the complexity
of re-running the full simulation pipeline for each trajectory.

\subsection{Sinusoidal Drift Trajectory}
\label{sec:trajectory_gen}

Rather than a random walk, which can accumulate arbitrarily large displacements,
the trajectory is modelled as an independent sinusoid on each spatial axis:
%
\begin{equation}
  \tau_a(t) = A_a \sin\!\left(2\pi f_a\,t + \varphi_a\right), \qquad a \in \{y,\,x\}.
  \label{eq:sinusoid}
\end{equation}
%
The normalised frequency is $f_a = n_a / T$, where
$n_a \sim \mathrm{Uniform}(1.5,\,3.5)$ is the number of oscillation cycles over the
$T$-frame sequence, and $\varphi_a \sim \mathrm{Uniform}(0,\,2\pi)$ is a random phase
that breaks any correlation between axes.
The amplitude $A_a$ is derived from a user-specified peak velocity $v$
(in LR pixels per frame): since the instantaneous peak velocity of a sinusoid is
$2\pi f_a A_a$, setting
%
\begin{equation}
  A_a = \min\!\left(\frac{v}{2\pi f_a},\;d_{\max}\right)
  \label{eq:amplitude}
\end{equation}
%
guarantees that the maximum drift speed equals exactly $v$, with $d_{\max}$ as a
hard bound that prevents large zero-padded border regions.
This parameterisation decouples biological realism, controlled by $v$, from
trajectory shape, controlled by $n_a$ and $\varphi_a$.
The sinusoidal model also naturally captures the direction reversals that
are characteristic of organelle motion constrained by cytoskeletal tracks.

\begin{algorithm}[h]
\caption{Sinusoidal drift trajectory generation}
\label{alg:trajectory}
\begin{algorithmic}[1]
  \Require $T$ (frames), $v$ (peak velocity in LR\,px/frame),
           $d_{\max}$ (max displacement per axis)
  \For{each axis $a \in \{y,\,x\}$}
    \State Draw $n_a \sim \mathrm{Uniform}(1.5,\,3.5)$;\quad set $f_a \leftarrow n_a / T$
    \State Draw $\varphi_a \sim \mathrm{Uniform}(0,\,2\pi)$
    \State $A_a \leftarrow \min\!\left(v\,/\,(2\pi f_a),\;d_{\max}\right)$
    \State $\boldsymbol{\tau}_{:,a} \leftarrow A_a\,\sin(2\pi f_a\,\mathbf{t} + \varphi_a)$,
           \quad $\mathbf{t} = [0,\,1,\,\ldots,\,T{-}1]^\top$
  \EndFor
  \State \Return $\boldsymbol{\tau} \in \mathbb{R}^{T \times 2}$
\end{algorithmic}
\end{algorithm}

\subsection{Frame Shifting and Per-Frame Targets}
\label{sec:frame_shifting}

Given a trajectory $\boldsymbol{\tau}$, each LR frame $\mathbf{L}[t]$ is shifted by
$(dy_t, dx_t) = \boldsymbol{\tau}[t]$ using bilinear interpolation with zero-padding.
The corresponding HR target is shifted by the same displacement scaled by the
upsampling ratio $s = H_\text{hr}/H = 121/64 \approx 1.89$, consistent with the
$2n-7$ output size.

\begin{algorithm}[h]
\caption{Per-frame shifting for motion augmentation}
\label{alg:shift}
\begin{algorithmic}[1]
  \Require $\mathbf{L}$ ($T \times H \times W$), $\mathbf{Y}$ ($H_\text{hr} \times W_\text{hr}$),
           $\boldsymbol{\tau}$ ($T \times 2$)
  \State $s_y \leftarrow H_\text{hr} / H$, \quad $s_x \leftarrow W_\text{hr} / W$
  \For{$t = 1, \ldots, T$}
    \State $(dy_t, dx_t) \leftarrow \boldsymbol{\tau}[t]$
    \State $\hat{\mathbf{l}}_t \leftarrow \mathrm{shift}(\mathbf{L}[t],\;(dy_t,\,dx_t))$
           \hfill\textit{// bilinear, zero-pad}
    \State $\hat{\mathbf{y}}_t \leftarrow \mathrm{shift}(\mathbf{Y},\;(dy_t s_y,\,dx_t s_x))$
           \hfill\textit{// bilinear, zero-pad}
  \EndFor
  \State \Return $\hat{\mathbf{L}}$ ($T \times H \times W$),\;
                 $\hat{\mathbf{Y}}$ ($T \times H_\text{hr} \times W_\text{hr}$)
\end{algorithmic}
\end{algorithm}

The critical consequence is that each frame $t$ receives its own HR ground-truth
$\hat{\mathbf{y}}_t$.
In the static case, a single HR image is broadcast across all time steps, so the
model can ignore temporal position entirely.
With motion augmentation, the loss $\mathcal{L}(\hat{Y}_t,\,\hat{\mathbf{y}}_t)$
forces the streaming model to predict the HR image at the \emph{current} position
of the field of view rather than a single consensus reconstruction.

\subsection{Relevance to Physical Speed}
\label{sec:speed_calibration}

The peak velocity $v$ in LR pixels per frame maps to a physical drift speed via
%
\begin{equation}
  v_\text{real}\;[\mu\text{m/s}]
    \;=\; v\;[\text{px/frame}]
    \;\times\; r\;[\text{frame/s}]
    \;\times\; p\;[\mu\text{m/px}],
  \label{eq:speed_calibration}
\end{equation}
%
where $r = 100$\,Hz is the camera acquisition rate and $p = 0.1$\,\textmu m/px
(100\,nm per pixel, consistent with the simulated dataset described in
Section~\ref{sec:synthetic}).
The configured speed range $[v_{\min},\,v_{\max}] = [0.2,\,1.5]$\,LR\,px/frame
therefore spans
%
\begin{align*}
  v_{\min} &= 0.2 \times 100 \times 0.1 = 1.0\; \textmu \text{m/s}, \\
  v_{\max} &= 1.5 \times 50 \times 0.1 = 7.5\;  \text{\textmu m/s}.
\end{align*}
%
This range covers the full biological drift envelope reported for vesicles and
mitochondria (0.5--5\,\textmu m/s~\cite{tang2013fast}), with 7.5\,\textmu m/s
providing headroom for fast kinesin-driven transport.
Both bounds are exposed as configuration parameters (\texttt{speed\_min},
\texttt{speed\_max}), making it straightforward to restrict training to slow
passive diffusion or to extend coverage to faster motor-driven motion.

\subsection{Integration into Training}
\label{sec:motion_training_integration}

Motion augmentation is applied on-the-fly inside the
\texttt{StreamingTensorDataset} data loader.
For each training sample, a motion augmentation is applied with a probability of $p_\text{motion} = 0.7$
This way model remembers to reconstruct static samples.
When augmentation is applied, peak velocity is drawn uniformly from $[0.2,\,1.5]$\,LR\,px/frame,
a trajectory is generated via Algorithm~\ref{alg:trajectory}, and the frames and
targets are shifted via Algorithm~\ref{alg:shift} with $d_{\max} = 38$\,LR\,px.
Otherwise the static HR image is kept across all $T$ frames.